\documentclass[11pt,letterpaper]{article}
\usepackage[spanish,es-nodecimaldot,es-noquoting]{babel}
\usepackage[T1]{fontenc}
\usepackage{lmodern,microtype,amsmath,amssymb,graphicx,booktabs}
\usepackage[margin=0.92in]{geometry}
\usepackage[authoryear,round]{natbib}
\usepackage{float}
\usepackage[hidelinks]{hyperref}
\hypersetup{pdftitle={Violencia criminal y accidentes registrados: rediseño causal},pdfauthor={Proyecto de Inferencia Causal, ITAM}}
\urlstyle{same}
\setlength{\parindent}{1.1em}
\setlength{\parskip}{0.15em}
\setlength{\emergencystretch}{2em}
\setlength{\tabcolsep}{3pt}
\renewcommand{\arraystretch}{1.12}
\linespread{1.06}
\setlength{\bibsep}{0pt}
\graphicspath{{../../data/outputs/figuras/}}
\input{../../data/outputs/tablas/causal_cifras.tex}
\newcommand{\tabla}[1]{\input{../../data/outputs/tablas/#1.tex}}
\title{\textbf{Violencia criminal y accidentes registrados}\\[0.2em]\large Jornadas extensas en municipios mexicanos, 2019--2025}
\author{Proyecto de Inferencia Causal · ITAM}\date{Noviembre de 2026}
\begin{document}\maketitle
\begin{abstract}
Estudio el cambio inmediato en accidentes registrados tras el inicio local de
jornadas extensas de violencia. El rediseño utiliza tres semanas previas del
mismo día de semana, emparejamiento individual y elegibilidad determinada con
información previa. En \RTrat{} pares tratados y \RFechas{} fechas, el
contraste de dos días es \RMainEfecto{} en $\log(1+Y)$, con IC de dos vías
[\RDobleIcinf{}; \RDobleIcsup{}]. El placebo no reproduce ese descenso.
Los diagnósticos previos y la sensibilidad a escala, calendario e inferencia
impiden una conclusión causal incondicional. Distingo este resultado de la
persistencia seleccionada posteriormente y de Culiacán descriptivo.
\end{abstract}
\section{Contexto y motivación}
¿Cómo cambian los accidentes registrados durante los dos primeros días de una
jornada extensa de violencia y qué ocurre después? ATUS registró 380\,391
accidentes en 2025 en vías urbanas y suburbanas no federales \citep{inegiatus}.
Los bloqueos, la suspensión del transporte y el retiro de personas de las calles
pueden modificar rápidamente la circulación. Los ataques a transportistas de
Chilpancingo en febrero de 2024 ilustran una suspensión del servicio
\citep{amapola}. También pueden cambiar las rutas, el riesgo por viaje y la
capacidad de captar siniestros. El resultado de interés es el registro publicado,
no una tasa de riesgo vial por viaje.

El objetivo es el ATT sobre pares municipio--episodio con elegibilidad y soporte
previos: el efecto de un inicio frente a su ausencia, manteniendo el historial
anterior. Cada par tratado recibe el mismo peso. El contraste principal promedia
las fechas 0 y 1; las fechas 2--6 y 7--13 examinan la trayectoria posterior.
Son dos días naturales, no exactamente 48 horas transcurridas desde un inicio
vespertino. El día 0 puede incluir horas sin exposición y el día 1 es el primer
día natural completo posterior.

\citet{robles} estudian consecuencias económicas de la violencia y \citet{dell}
documenta desvíos de rutas que motivan cautela sobre interferencia.
\citet{bauern} encuentran cambios viales contemporáneos durante campañas
alemanas de un día y distinguen ausencia de persistencia detectable de un efecto
permanente. Su exclusión de feriados y fines de semana largos motiva aquí
sensibilidades de calendario, sin equiparar ambas intervenciones.

Menos movilidad predice una caída, pero no determina el signo: los desvíos pueden
elevar el riesgo y una interrupción del registro puede reducir el conteo sin
mejorar la seguridad. La interpretación causal requiere supuestos sobre cambios
contrafactuales y captación. Este rediseño es posterior a observar T05; conserva
la referencia anterior en el apéndice y no se presenta como preespecificación.

\clearpage\section{Datos y muestra}
\paragraph{Accidentes y exposición.}
Uso el ZIP definitivo ATUS 1997--2025 del INEGI, fijado por SHA-256; analizo
2019--2025 y uso finales de 2018 para la historia de reporte
\citep{inegiatus,inegimeta}. Se conservan 2\,518\,436 registros de fecha válida
en 2019--2025. Cinco filas de fecha inválida se documentan sin redistribuirlas.
Sin identificador de accidente, se retienen atributos repetidos: la llave técnica
es año y renglón del CSV. Hora y gravedad concilian con el total diario.

Una jornada extensa reúne acciones coordinadas en tres o más lugares en 24
horas con afectación del espacio público, comercios o transporte. El catálogo
documenta fuentes, fechas locales y orígenes independientes; una jornada que
cruza medianoche puede tener fechas municipales distintas. Los incidentes menores
o inciertos delimitan antecedentes. No es un censo exhaustivo de violencia.
Chilchota (14/07/2021) y Contepec (23/04/2025) mantienen corroboración pendiente
y tienen una sensibilidad específica.

\paragraph{Cobertura y ceros.}
La muestra principal exige registros en cada uno de los tres meses completos
anteriores al mes de $t$, consultando anomalías de esos meses. Una anomalía sin
mes del año corriente no se supone anterior al inicio. Ningún resultado,
continuidad mensual, réplica o aislamiento posterior determina elegibilidad.
El resultado cuenta filas de fecha válida del release: un día sin filas es cero
\emph{registros publicados} en ese universo previamente definido, no ausencia
comprobada de accidentes. No se rellenan los faltantes del panel de continuidad.
La posible interrupción de captación posterior limita la interpretación.

\paragraph{Covariables y territorio.}
Población y viviendas proceden de EIC 2015 e ITER 2020; empleo, de SAIC
(Censos Económicos 2018 y 2023); autos, de VMRC 2018--2024; denominadores
anuales, de CONAPO; AGEEML evalúa equivalencias y cambios territoriales
\citep{eic2015,iter2020,saic,vmrc,conapo,ageeml}. Se selecciona la última
medición cerrada antes de $t-28$, conservando el cierre especial de Colima en
VMRC y de sus denominadores. Cero, confidencialidad, clave ausente y territorio
incompatible se distinguen. Son ediciones retrospectivas, no un archivo de
versiones disponibles en tiempo real en cada inicio.

Antes del matching, las uniones de controles son \RUnionCensoC{} (EIC/ITER),
\RUnionSaicC{} (SAIC) y \RUnionVmrcC{} (VMRC); para tratados son
\RUnionCensoT{}, \RUnionSaicT{} y \RUnionVmrcT{}, respectivamente. El
apéndice distingue unión y valor utilizable y reporta CONAPO y AGEEML.
\begin{table}[H]\centering\small
\caption{Comparabilidad previa de la muestra final}\label{tab:balance}
\tabla{causal_balance}
\par\smallskip\footnotesize DE ponderada entre pares. DME utiliza como escala
común la raíz del promedio de varianzas de ambos grupos de la referencia con
calibre individual. El nivel previo es la media de $[-28,-15]$; su DE entre
pares difiere de la DE de niveles diarios que estandariza los efectos.
\end{table}

\clearpage\section{Estrategia empírica}
\paragraph{Estimando y base.}
Sea $t$ el inicio local, $D_{it}$ la asignación focal y $A_{it}$ la elegibilidad
previa. Para $g(y)=\log(1+y)$, principal, o $g(y)=y$, complementaria, el objetivo es
\[
\tau_h^g=E[g(Y_{i,t+h}(1))-g(Y_{i,t+h}(0))\mid D_{it}=1,A_{it}=1,\text{soporte}].
\]
El régimen con inicio permite su trayectoria posterior: no condiciono el efecto
inmediato a que no cause episodios nuevos. Defino $q(h)=h\bmod7\in\{0,\ldots,6\}$,
$r(t,h)=t+q(h)$ y
\[
B_{ith}^g=\tfrac13\sum_{k\in\{28,21,14\}}g(Y_{i,r(t,h)-k}),\qquad
\Delta_{ith}^g=g(Y_{i,t+h})-B_{ith}^g.
\]
Las bases están en $[t-28,t-8]$ y son del mismo día de semana que el resultado.
Para $h=1$ son $\{-27,-20,-13\}$ y para $h=7$, $\{-28,-21,-14\}$.
Usar literalmente $r=t+h$ saldría de la ventana para ciertos horizontes.
Cada tramo promedia sus contrastes diarios. En log se transforma antes de
promediar y se admiten ceros. Como en Lecture 5, el log modifica paralelismo e
influencia de conteos altos; no garantiza identificación ni invariancia a escala.

\paragraph{Comparadores y pesos.}
Se excluyen antecedentes de exposición en $[t-49,t-1]$, incluidos 21 días de
lavado anteriores a la primera medición. Los controles están fuera de los
estados focales de la cohorte y de pistas estatales previas. Las nuevas
exposiciones posteriores no seleccionan el principal. Los cambios territoriales
se evalúan hasta $t-1$. La extinción del arrastre sigue siendo un supuesto.

Para cada tratado, $(pre\_nivel_c+0.5)/(pre\_nivel_i+0.5)$ debe estar entre
0.5 y 2. Después, población, empleo por mil y autos por mil deben quedar entre
la mitad y el doble de su propio valor. Se exigen diez candidatos de cinco
estados por tratado. En 1:muchos se reparte un peso entre sus donantes; en 1:1
se elige el vecino de menor distancia euclidiana en diferencias $\log_2$
relativas, con reemplazo y desempate por clave. El soporte se evalúa antes de
elegir ese único vecino. Los controles compartidos acumulan contribuciones.

Se estima por mínimos cuadrados ponderados con efectos fijos de fecha:
\[
\Delta_{ith}^g=\lambda_{th}+\beta_h^gD_{it}+X_{it}'\gamma_h+u_{ith}.
\]
M0 omite $X$: tratados pesan uno y controles tienen igual masa total por fecha.
La pendiente promedia por igual los DiD de cada tratado. M1 agrega log población;
M2, empleo por mil; M3, autos por mil. Muestra y pesos se mantienen en cada
progresión. El ajuste lineal puede modificar la ponderación implícita del ATT.

\paragraph{Identificación y sesgo.}
Sin anticipación en la base, el contraste incorpora $\tau_h^g$ más el cambio
de brecha contrafactual $B_h^g-\overline B_{\mathrm{base},h}^g$. El paralelismo
condicional elimina ese término, no una igualdad de niveles. Paralelismo,
ausencia de arrastre e interferencia y estabilidad de medición son supuestos
\emph{no comprobables directamente}. Cronologías, soporte, pesos y conciliación
son comprobables; previos y placebo son diagnósticos. Un pretest no significativo
no demuestra el supuesto \citep{roth}. El tratamiento recurrente no satisface
automáticamente la adopción permanente de otros estimadores \citep{dube}.

\clearpage\subsection{Inferencia y separación de estimandos}
La inferencia principal es CRV1 por estado, permitiendo dependencia regional y
serial. Usa $G/(G-1)(N-1)/(N-K)$ y $G-1$ grados de libertad; $K$ incluye
las fechas. La doble suma los componentes de estado y fecha y resta su intersección,
usando el mínimo de conglomerados para corrección y grados de libertad. No
cubre cualquier dependencia entre estados en fechas distintas. Las varianzas
indefinidas se identifican sin fabricar intervalos.

Hay \REstados{} estados, \REstadosTrat{} tratados y \REstadosEf{} estados
tratados de peso equivalente. CRV3 y las omisiones estatales conservan muestra
y pesos, sin rematching ni renormalización. Cada réplica vuelve a estimar fechas;
la varianza es $(G-1)/G\sum_g(\hat\beta_{-g}-\hat\beta)^2$. No se confunde
con el jackknife centrado en la media de réplicas \citep{mackinnon,jackknife}.

La rama de persistencia exige además continuidad observada en $[-28,+13]$ y
ausencia de otra exposición hasta $+14$, salvo el episodio focal. Su inclusión
puede cambiar bajo el tratamiento: no hereda automáticamente la interpretación
causal principal. También estimo 2--6 y 7--13 en la muestra previa para hacer
visible la selección. Son preguntas y poblaciones distintas.

\section{Resultados}
\paragraph{Efecto inmediato y escala.}
El principal contiene \RTrat{} tratados, \RCont{} controles municipio--fecha,
\RFechas{} fechas y \REpis{} episodios. \RLecturaPrincipal{}
El coeficiente log corresponde a \RPorcentajeUnoMas{}\,\% en la escala
geométrica de $1+Y$, no automáticamente en el promedio aritmético de $Y$.
En niveles, el cambio es \RNivelEfecto{} registros por día y municipio,
\RNivelEfectosigma{} DE previas, con IC doble [\RNivelDobleIcinf{};
\RNivelDobleIcsup{}]. Las escalas no aportan igual evidencia estadística.
\RLecturaCRV{}
\begin{table}[H]\centering\small
\caption{Contrastes principales, matching 1:muchos}\label{tab:main}
\tabla{causal_principales}
\par\smallskip\footnotesize EE estatal entre paréntesis; IC doble estado--fecha.
M0, $N=\RN{}$ pares. Los tramos son promedios diarios. La DE diaria se calcula
en $[-28,-8]$, \texttt{ddof=1}; en niveles es \RSigmaPre{}.
\end{table}

El MDE supone prueba bilateral al 5\,\%, potencia 80\,\%, EE fijo e inversión
de $t$ no central. En niveles es \RNivelMde{} por estado y
\RNivelDobleMde{} en dos vías: \RMdePctNivel{}\,\% y
\RMdePctNivelDoble{}\,\% de la media previa \RMediaPre{}.
En log es \RMainMde{} y \RDobleMde{}. Describe potencia condicional,
no corrige sesgo ni demuestra equivalencia con cero. El apéndice presenta
1:1, DME y M0--M3 con sus muestras constantes.

\clearpage\paragraph{Trayectoria y falsificaciones.}
La figura \ref{fig:final} muestra ambas escalas y matchings. La semana previa
log es \RPreviaEfecto{}, con IC doble [\RPreviaDobleIcinf{};
\RPreviaDobleIcsup{}]. Persisten caídas puntuales en $h=-2$
(\RMenosDosEfecto{}) y $h=-5$ (\RMenosCincoEfecto{}). El agregado puede
esconder desviaciones diarias; su falta de significancia no valida tendencias
paralelas. Los intervalos de la figura son puntuales.
\begin{figure}[H]\centering
\includegraphics[width=\linewidth]{causal_dinamica.pdf}
\caption{Cambios respecto a tres semanas limpias del mismo día de semana.
Muestra determinada antes del inicio; IC puntuales de dos vías.}\label{fig:final}
\end{figure}

El Wald compara semanas -4 ($[-28,-22]$), -3 ($[-21,-15]$) y -1
($[-7,-1]$), omitiendo -2 ($[-14,-8]$). Es una normalización diagnóstica
distinta de la base principal. Evita la dependencia exacta entre las tres
semanas de la base normalizadas por su propio promedio. Su $p$ es \RWaldP{}
por estado y \RWaldDobleP{} en dos vías. El rechazo estatal debilita el
paralelismo: el rediseño no elimina todas las objeciones de identificación.
Se publican rango, autovalor mínimo y condicionamiento.

El placebo mueve el inicio a $t'=t-28$, reconstruyendo elegibilidad y matching;
sus bases son $[t-56,t-36]$ respecto al inicio real. En d0..1, log es
\RPlaceboEfecto{}, con IC doble [\RPlaceboIcinf{}; \RPlaceboIcsup{}].
CRV1 estatal sí rechaza 2--6 en ambas escalas y semana 0 en niveles; dos vías
no rechaza. No es un placebo uniformemente nulo. El apéndice presenta seis
horizontes con muestras propias; no descarta todos los confusores de las fechas reales.

\paragraph{Persistencia y sensibilidades.}
\RLecturaPersistencia{} La señal en semana +1 depende de inferencia y
exclusiones: no afirmo disipación completa desde el día 2. Con base $wk-1$,
log d0..1 es \RBaseSemanaEfecto{} [\RBaseSemanaIcinf{}; \RBaseSemanaIcsup{}].
El apéndice informa todas las sensibilidades y sus muestras.

\clearpage\paragraph{Calendario y controles.}
Las cinco fechas propuestas no son feriados laborales generales LFT: son
recesos de alumnos de educación básica, con talleres y el domingo previo al
regreso de enero de 2024. También hay episodios durante verano y otros recesos.
Conservo la comparación de cinco fechas y añado excluir cohortes cuyo día 0 o
1 coincida con feriado LFT o receso/suspensión SEP \citep{lft,calendariosep}.
Los fines de semana ordinarios no se excluyen por sí solos; los que integran
un receso continuo se incluyen hasta el regreso. El receso extraordinario de
2020 se distingue de vacaciones ordinarias. La aproximación nacional de
educación básica no cubre todas las adaptaciones estatales o universitarias.
La exclusión amplia reduce el soporte y su intervalo log inmediato incluye cero.

El calibre anterior dejaba Morelia con \RMoreliaAntes{} donantes el
23/04/2025. El universo previo final ofrece \RMoreliaUniverso{} controles;
tras el calibre individual quedan \RMoreliaCalibre{} y después de las tres
covariables, \RMoreliaDespues{}. Al excluir señales conserva
\RMoreliaSenales{}. Así se distingue depuración de controles de la antigua
pérdida composicional de Morelia. Se publican las listas, sin imponer vecinos.

\paragraph{Cruce exploratorio de hora y gravedad.}
El cuadro \ref{tab:cross} usa la misma muestra y base limpia, en niveles.
El contraste pareado conserva covarianza entre franjas. Las nueve celdas,
incluidas horas desconocidas, concilian con el total. Esta exploración de varios
resultados no identifica mecanismos separados y no corrige pruebas múltiples.
\begin{table}[H]\centering\small
\caption{d0..1: noche menos día por gravedad}\label{tab:cross}
\tabla{causal_cruce}
\par\smallskip\footnotesize Registros por municipio y día; IC de dos vías.
\end{table}

\paragraph{Culiacán, descriptivo.}
La serie comienza el 9/09/2024 \citep{luz,sinaloa}. Se suman Culiacán y
Eldorado desde noviembre sólo si ambas claves tienen continuidad \citep{decreto}.
Se retiran donantes con cambios tipo M: San Luis Potosí 24028 se divide con
Villa de Pozos 24059 el 30/08/2024. Se eligen \RCulDonantes{} donantes y
la brecha media es \RCulBrecha{} por semana, en \RCulObs{} de \RCulTotal{}
semanas observables. ATUS no tiene filas del 7 al 27 de octubre de 2024; la
continuidad mensual infiere ceros. La reanudación de tránsito y otra fuente de
distinto universo señalan problemas de captación \citep{transitocul,mapasin}.
La brecha no identifica el efecto causal de la crisis.

\section{Conclusión y validez externa}
El rediseño es compatible con una reducción inmediata de registros,
especialmente en log, que el placebo no reproduce. La atribución causal sigue
condicionada a paralelismo, interferencia y medición. Los previos estatales,
$h=-2$, escala, calendario e indicios de semana +1 impiden defender sin reservas
una reducción causal que desaparezca por completo desde el día 2. Persistencia
seleccionada y Culiacán responden a preguntas distintas. El alcance se limita
a jornadas documentadas con soporte; no se extrapola automáticamente a violencia
menor, vías federales o riesgo por viaje. Separar movilidad, riesgo y captación
requiere aforos y evidencia de registro adicionales.

\clearpage\begingroup\small\raggedright
\renewcommand{\url}[1]{\href{#1}{Fuente}}
\bibliographystyle{plainnat}\bibliography{referencias}\endgroup

\clearpage\appendix\section{Fuentes y especificaciones progresivas}
\begin{table}[H]\centering\small\caption{Uniones antes del matching}
\tabla{causal_uniones}\end{table}
El universo son pares elegibles antes de exigir valores completos. Unión no
equivale a valor utilizable. CONAPO corresponde al año del numerador. AGEEML
es un inventario de cambios: no tener fila no es fallo de unión; se evalúa
estabilidad territorial entre fechas.
\begin{table}[H]\centering\small\caption{Uniones auxiliares y territorio}
\tabla{causal_uniones_aux}\end{table}
\begin{table}[H]\centering\small\caption{Pérdidas de tratados hasta la muestra principal}
\tabla{causal_perdidas}\end{table}
Cada fila cuenta pérdidas después de los filtros anteriores. Un fallo de soporte
significa menos de diez candidatos o menos de cinco estados. La elegibilidad
se construye hasta $t-1$; el soporte usa sólo covariables y niveles previos.
Las exclusiones de persistencia se contabilizan aparte y no seleccionan el
efecto inmediato principal. Los faltantes de covariables no se convierten en cero.

\clearpage
\begin{table}[H]\centering\footnotesize\caption{d0..1: M0--M3, muestra fija}
\tabla{causal_modelos}\end{table}
M0: tratamiento y fecha; M1: añade log población; M2: empleo/mil; M3: autos/mil.
La muestra completa exige covariables utilizables y calibre individual, sin
los tres intervalos adicionales del matching; conserva todos los tratados con
ese soporte. Cada progresión usa exactamente sus mismos pares y pesos.
La referencia sin matching mantiene los mismos tratados y controles completos,
con pesos uniformes por fecha. Los matchings comparten tratados, no donantes.
Las DME del cuadro 1 usan una escala común fijada antes del matching.
\begin{table}[H]\centering\small
\caption{DME con los mismos tratados antes y después del matching}
\tabla{causal_balance_prepost}\end{table}
La referencia previa incluye controles con datos completos y ponderación
uniforme dentro de cada fecha. Las columnas posteriores cambian los vínculos
y pesos manteniendo tratados. La escala de estandarización es común, fijada
en la referencia con calibre individual; una DME pequeña no demuestra
balance de variables omitidas ni paralelismo. En niveles y log las covariables
son las mismas, por lo que no se duplican estas estadísticas de balance.

\clearpage\section{Persistencia y placebo}
\begin{table}[H]\centering\footnotesize\caption{Log: selección y horizontes, IC doble}
\tabla{causal_persistencia}\end{table}
La rama de aislamiento exige cobertura $[-28,+13]$ y ausencia de otra exposición
hasta +14. Sus diferencias incorporan selección posterior. El placebo reconstruye
la muestra en $t-28$; sus tamaños no tienen que coincidir con los reales.
\begin{table}[H]\centering\small\caption{Placebo -28: seis horizontes solicitados}
\tabla{causal_placebo}\end{table}
$p_E$ corresponde a CRV1 estatal; los intervalos son de dos vías. Los rechazos
estatales en tramos posteriores impiden describir todos los placebos como ceros limpios.
Se exportan inicio real, índice ficticio, resultado y las tres fechas de base en
\path{data/outputs/modelos/causal_fechas_base.csv}. También se publican días
individuales -7 a -1 y 0 a 13 y la sensibilidad con base $wk-1$ en exactamente
la misma muestra final. Es un ejercicio posterior a T05, no una referencia
elegida después para seleccionar significancia.

\clearpage\section{Sensibilidades y calendario}
\begin{table}[H]\centering\footnotesize\caption{d0..1: exclusiones, IC doble}
\tabla{causal_sensibilidades}\end{table}
Las cinco fechas son 31/12/2019, 05/01/2023, 09/04/2023, 07/01/2024 y
23/04/2025. La regla amplia se aplica al día 0 o 1 sin eliminar automáticamente
domingos. Se distinguen vacaciones, talleres, suspensiones y receso extraordinario
de 2020; no se llama vacaciones a toda la educación a distancia.

Se excluyen por separado Chilchota 16025 (14/07/2021) y Contepec 16017
(23/04/2025). Se añade excluir los cuatro pendientes de auditoría: también
Huaniqueo 16037 (23/04/2025) y San Fernando 28035 (30/04/2023). Este último
entra en el rediseño; Huaniqueo sigue fuera por elegibilidad previa.
Las señales son Centro 27004 (01--02/04/2025) y Coatzacoalcos
30039 (20/08/2022). Se elimina su contribución si intersecta $[t-49,t+14]$ y
se reconstruyen matching y pesos. Coatzacoalcos ocurrió después del 09/08/2022:
es una sensibilidad ex post explícita, no un filtro del principal.
El cotejo anterior cubrió 27 pares de cuatro fechas. ACLED devolvió HTTP 403:
no existe verificación exhaustiva y una búsqueda sin señal no demuestra ausencia.
\begin{table}[H]\centering\footnotesize\caption{Soporte de los escenarios}
\tabla{causal_soporte}\end{table}

\clearpage\begin{table}[H]\centering\small
\caption{Cohortes finales y calendario oficial}\tabla{causal_calendario}\end{table}
El calendario es nacional de educación básica, con fines de semana contiguos
hasta la reanudación. No representa cada institución o adaptación estatal.
La LFT sigue el artículo 74, incluido el cambio presidencial de 2024. Ninguna
cohorte estimada coincide en d0..1 con una jornada electoral ordinaria.

\clearpage\section{Inferencia y potencia}
\begin{table}[H]\centering\small\caption{d0..1: agrupamiento, CRV3 y MDE}
\tabla{causal_inferencia}\end{table}
MDE bilateral al 5\,\%, potencia 80\,\%. La DE previa usa valores diarios en
$[-28,-8]$, \texttt{ddof=1}; la descriptiva usa dispersión entre medias de
$[-28,-15]$. Un MDE log no es un porcentaje directo de la media de $Y$.
\begin{table}[H]\centering\footnotesize\caption{Wald semanal de tres restricciones}
\tabla{causal_previos}\end{table}
Se omite semana -2 y se contrastan -4, -3 y -1. Se informa singularidad o
indefinición sin usar una inversa generalizada para fabricar $p$. CRV3 tampoco
repara anticipación, selección o interferencia.

\clearpage\begin{table}[H]\centering\small
\caption{Omisión de estados, M0 y 1:muchos}\tabla{causal_omision}\end{table}
Las réplicas mantienen pesos originales y vuelven a absorber las fechas. No se
fuerza igualdad de masas después de omitir un estado: renormalizar o rematchear
sería otro estimador, no el CRV3 anunciado. Se exportan 1:1 y M1--M3.

\clearpage\section{T05 y trazabilidad del rediseño}
T05 se fijó el 17/11/2026 antes de observar efectos: base -2, calibre por media
de fecha, continuidad $[-28,+15]$, lavado 21 días y aislamiento hasta +16.
Su manifiesto registra 77 tratados, 7\,796 controles y 31 fechas. La revisión
de inventario produjo 73 tratados, 7\,351 controles y 30 fechas. Se reconstruye
esa revisión como referencia separada, no como el catálogo original de 77.
T05 y sus huellas permanecen intactos. La comparación utiliza la referencia
revisada identificada, sin confundirla con el catálogo original.

El nuevo diseño cambia base, escala, elegibilidad y calibre después de observar
T05. No incorpora $wk-1$ a la base limpia. La auditoría mezclaba tres semanas
con algunos ejercicios de dos; aquí el principal utiliza tres. Las dos escalas
y todas las sensibilidades se mantienen independientemente de su significancia.
El nuevo número de 77 tratados coincide con T05, pero sus pares, fechas y pesos
no son idénticos.
\begin{table}[H]\centering\footnotesize
\caption{Cambios de muestra y base separados; IC estatales}\tabla{causal_trazabilidad}\end{table}
La referencia conserva la selección auditada posterior. Las dos últimas columnas
conservan exactamente la muestra final y sólo cambian la base. El código y CSV
identifican pertenencia y pesos. La sensibilidad a $wk-1$ también mantiene esos
pares; no se recalcula soporte según sus resultados.
\begin{table}[H]\centering\footnotesize
\caption{d0..1 con idénticos pares y pesos: base e IC de dos vías}
\tabla{causal_bases}\end{table}

Morelia ilustra el mecanismo: el calibre antiguo la centraba en el promedio de
nueve tratados de la fecha. Quitar Centro la dejaba debajo de diez candidatos y
eliminaba al tratado. Por ello el cambio de efecto no medía solamente depuración
de controles. El calibre final se refiere a su propio nivel. Las listas completas
se publican en \path{data/outputs/modelos/causal_morelia_donantes.csv}.

Se preservan catálogo y fuentes auditadas; los pendientes tienen sensibilidades,
no recodificación oportunista. El manifiesto histórico sigue verificando datos y
T05 y su muestra esperada. Sus hashes de código describen la versión anterior;
no se reemplazan para validar el código nuevo. El código y los productos finales
tienen evidencia SHA-256 separada. El cambio de muestra original a revisada se
documenta aquí sin citar un archivo ausente.

\clearpage\section{Culiacán y reproducción}
\begin{figure}[H]\centering\includegraphics[width=\linewidth]{culiacan.pdf}
\caption{Comparación semanal estrictamente descriptiva. Base -2; sombreado de
captación en revisión y huecos por cobertura insuficiente.}\end{figure}
El entrenamiento usa semanas -25 a -10 y validación -9 a -3, omitiendo de ella
la base -2. Se vuelven a elegir doce vecinos sin segregaciones M durante el
intervalo semanal relevante. Se preserva el empalme Culiacán--Eldorado desde
noviembre. Los cambios de nombre W no son divisiones. Se entregan serie,
donantes, descartes territoriales y placebos descriptivos; una semana sin cobertura
no se convierte en efecto causal.

El paquete \path{src/choques/}, configuración y pruebas acompañan al escrito.
\texttt{README.md} fija Python 3.13.2, lock y compilador. ATUS y los cien
originales municipales deben corresponder a sus manifiestos: una edición nueva
en la misma URL no los sustituye. Las entradas son
\path{data/inputs/episodios.csv} y \path{data/inputs/fuentes_eventos.csv},
reconstruidas desde las hojas de evidencia. Las rutas parten de la carpeta que
contiene \texttt{repo/}, \texttt{data/} y \texttt{tasks/}.

La CLI ejecuta \texttt{validar}, \texttt{todo} y \texttt{compilar}.
\texttt{python reproducir.py} crea otra copia, verifica originales, reconstruye
ocho etapas y exige igualdad SHA-256 de todos los productos y del PDF. Las
pruebas unitarias y sintéticas cubren calibre heterogéneo, base limpia, log con
ceros, invariancia al futuro, territorio, pesos e inferencia con refit de fechas.
Los tiempos variables quedan separados de productos deterministas. El registro
de verificación identifica comandos realmente ejecutados y revisión visual de
todas las páginas.

El profesor autorizó el uso de Python para este proyecto, según confirmación
del estudiante. T08 está fuera del alcance. Los datos públicos se enlazan en
la bibliografía y los manifiestos fijan sus ediciones.
\end{document}
